\subsection{Propósito, Alcance y Objetivos [Pendientes: 1. delimitar sector y cobertura, 2. incorporar la reformulación y decisiones pendientes del VFP]}\label{sec:spmp-6-2}
El proyecto busca concretar la visión anterior mediante el desarrollo de una aplicación de escritorio y de los medios necesarios para preparar, catalogar y consultar sus escenarios. El trabajo se limita a un sector de los Cerros Orientales de Bogotá todavía por delimitar y articula la preparación de datos, la configuración de WFDS, el precálculo mediante la infraestructura que se confirme y la integración del flujo principal de consulta. La distribución de las igniciones de referencia será de densidad variable guiada por amenaza de ignición, con mayor concentración en zonas de mayor propensión, y se mantendrá separada conceptualmente de la malla CFD que discretiza cada simulación. FARSITE pertenece a los antecedentes y queda excluido de la implementación, al igual que una aplicación web o móvil, la cobertura integral de los Cerros y las recomendaciones de reportes comunitarios o gestión de recursos que no se han incorporado al alcance~\cite{diseno2026,vfp2026}.

La formulación de trabajo del objetivo general consiste en desarrollar la aplicación para estimar y visualizar posibles comportamientos de incendios mediante escenarios WFDS precalculados, con el fin de apoyar al comandante o coordinador de la operación en el sector seleccionado. Los cuatro objetivos específicos se conservan como base del VFP y se organizan alrededor del mecanismo de búsqueda con condiciones meteorológicas, la construcción de la biblioteca, la configuración reproducible del simulador y la integración de la visualización en escritorio. Esta síntesis permite administrar el trabajo sin establecer una lista alternativa de objetivos aprobados. Cuando se acuerde su reformulación en el VFP, deberán revisarse conjuntamente el catálogo de entregables, la aceptación y la planificación, conservando en la trazabilidad la correspondencia con la formulación sustituida~\cite{vfp2026,pafp2026}.

El alcance temporal de este plan comprende la administración del desarrollo y de la entrega previstos para 2027-I, tomando como entrada el trabajo de preparación realizado o pendiente de Inicio. No se compromete una ampliación posterior del territorio ni un servicio de operación permanente. Las decisiones sobre horizontes de consulta, compatibilidad de escenarios, agregación, contingencia meteorológica y tamaño de la biblioteca deberán resolverse con evidencia y reflejarse en los requisitos, en lugar de fijarse mediante supuestos de este documento.
